\subsection*{一步一步建立直观}

\paragraph{1. 把每次增加一点理解为乘法更新}

实数复利每一步乘以 $1+r\Delta t$。若把增长率换成 $i$，一步变成乘以 $1+i\Delta t$，它近似一次小旋转，却还带着略大于一的模长。右图显示连续极限的圆周轨迹，而不是直接播放有限次复利步骤。



\paragraph{2. 为什么细分后多余增长会消失}

单步模长为 $\sqrt{1+\Delta t^2}$，偏离一的主要量是二阶；完成固定总时间需要的一阶数量级步数，不足以留下非零的二阶累计模长误差。另一方面，每步的一阶转角会积成有限总角度。这是旋转留下、额外膨胀消失的原因。



\paragraph{3. 真实径向增长与离散误差不同}

螺线情形故意加入非零实部，表示方程中真实存在的径向增长。它不会因为时间步细分而消失；有限步虚数复利的模长偏差则会在极限中消失。两幅外观相似的螺线可能来自不同机制，必须根据更新规则区分。
